\begin{abcliste}
\abc
With the angular frequency
\begin{align}
 \omega &= \omF = \frac{2\pi}{\T} = \om
\end{align}
the largest speed is
\begin{align}
 \hat v &= \vmaxF \\
   &= \A\times\om \\
   &= \vmax \approx \result{\vmaxP{0}{2}}
\end{align}

\abc
The kinetic energy grows with the square of the speed: at half the speed it is a quarter of the total energy $E$, so the potential energy (of the spring, measured from the resting point) is three quarters. Since $E_\mathrm{pot}/E = (y/A)^2$,
\begin{align}
 \left(\frac{y}{A}\right)^2 &= \frac{3}{4} \\
 y &= \yF \\
   &= \frac{\sqrt{3}}{2}\times\A \\
   &= \y \approx \result{\yP{-2}{2}}
\end{align}
Surprisingly far out: the baby still has half its top speed when it has covered \SI{87}{\percent} of the way to the turning point.
\end{abcliste}
